\clearpage
\lecturetitle{第18讲\quad 多元函数积分学（仅数学一）}

\parttitle{一}{三重积分}

\topictitle{1}{概念}

设 $f(x,y,z)$ 是空间有界闭区域 $\Omega$ 上的有界函数，将 $\Omega$ 任意分成 $n$ 个小闭区域
\[
\Delta v_1,\Delta v_2,\ldots,\Delta v_n,
\]
其中 $\Delta v_i$ 表示第 $i$ 个小闭区域，也表示它的体积。在每个 $\Delta v_i$ 上任取一点
$(\xi_i,\eta_i,\zeta_i)$，作乘积 $f(\xi_i,\eta_i,\zeta_i)\Delta v_i$，并作和
\[
\sum_{i=1}^{n}f(\xi_i,\eta_i,\zeta_i)\Delta v_i.
\]
如果当各小闭区域直径中的最大值 $\lambda$ 趋于零时，这个和的极限总存在，并且与 $\Delta v_i$ 的分法及 $(\xi_i,\eta_i,\zeta_i)$ 的取法均无关，则称此极限为函数 $f(x,y,z)$ 在闭区域 $\Omega$ 上的三重积分，记作
\[
\iiintO f(x,y,z)\dd v,
\]
即
\[
\iiintO f(x,y,z)\dd v
=\lim_{\lambda\to0}\sum_{i=1}^{n}f(\xi_i,\eta_i,\zeta_i)\Delta v_i.
\]
其中 $f(x,y,z)$ 称为被积函数，$f(x,y,z)\dd v$ 称为被积表达式，$\dd v$ 称为体积元素，$x,y,z$ 称为积分变量，$\Omega$ 称为积分区域。

若 $f(x,y,z)$ 在 $\Omega$ 上连续，则三重积分一定存在。

\begin{notebox}[三重积分的物理意义]
设一物体占有 $Oxyz$ 中的闭区域 $\Omega$，在点 $(x,y,z)$ 处的体密度为 $\rho(x,y,z)$，且 $\rho$ 在 $\Omega$ 上连续，则物体的质量为
\[
M=\iiintO\rho(x,y,z)\dd v.
\]
\end{notebox}

\topictitle{2}{性质}

以下总假设 $\Omega$ 为空间有界闭区域。

\prop{1（求空间区域的体积）}
\[
\iiintO1\dd v=\iiintO\dd v=V,
\]
其中 $V$ 为 $\Omega$ 的体积。

\prop{2（可积函数必有界）}
设 $f(x,y,z)$ 在 $\Omega$ 上可积，则其在 $\Omega$ 上必有界。

\prop{3（积分的线性性质）}
设 $k_1,k_2$ 为常数，则
\[
\iiintO[k_1f(x,y,z)\pm k_2g(x,y,z)]\dd v
=k_1\iiintO f(x,y,z)\dd v
\pm k_2\iiintO g(x,y,z)\dd v.
\]

\prop{4（积分的可加性）}
设 $f(x,y,z)$ 在 $\Omega$ 上可积，且 $\Omega_1\cup\Omega_2=\Omega$，$\Omega_1\cap\Omega_2=\varnothing$，则
\[
\iiintO f(x,y,z)\dd v
=\iiint_{\Omega_1}f(x,y,z)\dd v
+\iiint_{\Omega_2}f(x,y,z)\dd v.
\]

\prop{5（积分的保号性）}
设 $f(x,y,z),g(x,y,z)$ 在 $\Omega$ 上可积，且在 $\Omega$ 上 $f(x,y,z)\le g(x,y,z)$，则
\[
\iiintO f(x,y,z)\dd v\le\iiintO g(x,y,z)\dd v.
\]
特殊地，有
\[
\left|\iiintO f(x,y,z)\dd v\right|
\le\iiintO|f(x,y,z)|\dd v.
\]

\prop{6（三重积分的估值定理）}
设 $M,m$ 分别是 $f(x,y,z)$ 在 $\Omega$ 上的最大值和最小值，$V$ 为 $\Omega$ 的体积，则
\[
mV\le\iiintO f(x,y,z)\dd v\le MV.
\]

\prop{7（三重积分的中值定理）}
设 $f(x,y,z)$ 在 $\Omega$ 上连续，$V$ 为 $\Omega$ 的体积，则在 $\Omega$ 上至少存在一点 $(\xi,\eta,\zeta)$，使得
\[
\iiintO f(x,y,z)\dd v=f(\xi,\eta,\zeta)V.
\]

\topictitle{3}{普通对称性与轮换对称性}

（1）普通对称性。假设 $\Omega$ 关于 $xOz$ 面对称，$\Omega_1$ 是 $\Omega$ 在 $xOz$ 面右边的部分，则
\[
\iiintO f(x,y,z)\dd v=
\begin{cases}
2\displaystyle\iiint_{\Omega_1}f(x,y,z)\dd v,
&f(x,y,z)=f(x,-y,z),\\[2mm]
0,&f(x,y,z)=-f(x,-y,z).
\end{cases}
\]
关于其他坐标面对称的情况与此类似。

（2）轮换对称性。若把 $x$ 与 $y$ 对调后，$\Omega$ 不变，则
\[
\iiintO f(x,y,z)\dd v=\iiintO f(y,x,z)\dd v.
\]
关于其他变量的轮换情形与此类似。

\topictitle{4}{计算}

\paragraphline{（1）直角坐标系，$\dd v=\dd x\dd y\dd z$}

\circitem{1}
\key{先一后二法（先 $z$ 后 $xy$ 法，也叫投影穿线法）。}
若 $\Omega$ 在 $xOy$ 面上的投影区域为 $D_{xy}$，且
\[
z_1(x,y)\le z\le z_2(x,y),
\]
则
\[
\iiintO f(x,y,z)\dd v
=\iint_{D_{xy}}\dd x\dd y
\int_{z_1(x,y)}^{z_2(x,y)}f(x,y,z)\dd z.
\]

\begin{notebox}[定限原则]
后积先定限，限内画条线，先交写下限（下曲面），后交写上限（上曲面）。
\end{notebox}

\circitem{2}
\key{先二后一法（先 $xy$ 后 $z$ 法，也叫定限截面法）。}
若 $a\le z\le b$，以 $z$ 为常数截得的截面区域为 $D_z$，则
\[
\iiintO f(x,y,z)\dd v
=\int_a^b\dd z\iint_{D_z}f(x,y,z)\dd x\dd y.
\]

\paragraph{（2）柱面坐标系}
\[
x=r\cos\theta,\qquad y=r\sin\theta,\qquad z=z,
\]
\[
\dd x\dd y\dd z=r\dd r\dd\theta\dd z,
\]
从而
\[
\iiintO f(x,y,z)\dd x\dd y\dd z
=\iiint_{\Omega'}f(r\cos\theta,r\sin\theta,z)\,r\dd r\dd\theta\dd z.
\]

\paragraph{（3）球面坐标系}
\[
x=r\sin\varphi\cos\theta,\qquad
y=r\sin\varphi\sin\theta,\qquad
z=r\cos\varphi,
\]
\[
\dd x\dd y\dd z=r^2\sin\varphi\dd r\dd\varphi\dd\theta,
\]
从而
\[
\iiintO f(x,y,z)\dd x\dd y\dd z
=\iiint_{\Omega'}f(r\sin\varphi\cos\theta,r\sin\varphi\sin\theta,r\cos\varphi)
r^2\sin\varphi\dd r\dd\varphi\dd\theta.
\]

\begin{notebox}[球面坐标系下定限：转、张、穿]
第一步，\key{转}：从 $xOz$ 面出发，绕 $z$ 轴转动，确定 $\theta$；
第二步，\key{张}：从 $z$ 轴出发张开锥面，确定 $\varphi$；
第三步，\key{穿}：从原点发出射线穿过区域，确定 $r$。
\end{notebox}

\paragraph{一般换元法}

设
\[
x=x(u,v,w),\qquad y=y(u,v,w),\qquad z=z(u,v,w)
\]
是从空间 $(u,v,w)$ 到空间 $(x,y,z)$ 的一一映射，且 $x,y,z$ 具有一阶连续偏导数，并且
\[
\frac{\partial(x,y,z)}{\partial(u,v,w)}
=\begin{vmatrix}
x_u&x_v&x_w\\
y_u&y_v&y_w\\
z_u&z_v&z_w
\end{vmatrix}\ne0,
\]
则
\[
\iiintO f(x,y,z)\dd x\dd y\dd z
=\iiint_{\Omega'}f(x(u,v,w),y(u,v,w),z(u,v,w))
\left|\frac{\partial(x,y,z)}{\partial(u,v,w)}\right|
\dd u\dd v\dd w.
\]

\topictitle{5}{应用}

（1）若 $\Omega$ 是物体所占的空间区域，则物体的体积为 $V=\iiintO\dd v$。

（2）对于空间物体，若体密度为 $\rho(x,y,z)$，$\Omega$ 是物体所占的空间区域。

\noindent 物体的质量为 $M=\iiintO\rho(x,y,z)\dd v$，质心 $\bar x,\bar y,\bar z$ 为
\[
\bar x=\frac{\iiintO x\rho\dd v}{\iiintO\rho\dd v},\qquad
\bar y=\frac{\iiintO y\rho\dd v}{\iiintO\rho\dd v},\qquad
\bar z=\frac{\iiintO z\rho\dd v}{\iiintO\rho\dd v}.
\]

\begin{notebox}
（1）在考研范围内，重心就是质心。

（2）当密度 $\rho$ 为常数时，重心就成了形心。

（3）当密度 $\rho$ 为常数时，可反用形心公式
\[
\iiintO x\dd v=\bar xV,\qquad
\iiintO y\dd v=\bar yV,\qquad
\iiintO z\dd v=\bar zV.
\]
\end{notebox}

（3）计算物体对 $x$ 轴、$y$ 轴、$z$ 轴和原点 $O$ 的转动惯量：
\begin{align*}
I_x&=\iiintO(y^2+z^2)\rho(x,y,z)\dd v,&
I_y&=\iiintO(x^2+z^2)\rho(x,y,z)\dd v,\\
I_z&=\iiintO(x^2+y^2)\rho(x,y,z)\dd v,&
I_O&=\iiintO(x^2+y^2+z^2)\rho(x,y,z)\dd v.
\end{align*}

（4）物体对外部一点 $M_0(x_0,y_0,z_0)$ 处质量为 $m$ 的质点的引力。令
\[
r_0=\sqrt{(x-x_0)^2+(y-y_0)^2+(z-z_0)^2},
\]
则
\begin{align*}
F_x&=Gm\iiintO\frac{\rho(x,y,z)(x-x_0)}{r_0^3}\dd v,\\
F_y&=Gm\iiintO\frac{\rho(x,y,z)(y-y_0)}{r_0^3}\dd v,\\
F_z&=Gm\iiintO\frac{\rho(x,y,z)(z-z_0)}{r_0^3}\dd v.
\end{align*}

\parttitle{二}{第一型曲线积分}

\topictitle{1}{概念}

设 $f(x,y,z)$ 是光滑曲线弧 $\Gamma$ 上的有界函数，将 $\Gamma$ 任意分成 $n$ 个小弧段，弧长分别为 $\Delta s_i$，在第 $i$ 个小弧段上任取一点 $(\xi_i,\eta_i,\zeta_i)$。如果当各小弧段长度的最大值 $\lambda$ 趋于零时，极限
\[
\lim_{\lambda\to0}\sum_{i=1}^{n}f(\xi_i,\eta_i,\zeta_i)\Delta s_i
\]
总存在，并且与小弧段的分法及取点的方法均无关，则称此极限为函数 $f(x,y,z)$ 在曲线弧 $\Gamma$ 上对弧长的曲线积分或第一型曲线积分，记作
\[
\int_{\Gamma}f(x,y,z)\dd s
=\lim_{\lambda\to0}\sum_{i=1}^{n}f(\xi_i,\eta_i,\zeta_i)\Delta s_i.
\]

\begin{notebox}
若 $f(x,y,z)=\rho(x,y,z)$ 是曲线形细杆的线密度，则
\[
M=\int_{\Gamma}\rho(x,y,z)\dd s
\]
表示该细杆的质量。
\end{notebox}

\topictitle{2}{性质}

以下总假设 $\Gamma$ 为空间有限长分段光滑曲线。

\prop{1（求空间曲线的长度（弧长））}
\[
\int_{\Gamma}1\dd s=l_{\Gamma},
\]
其中 $l_{\Gamma}$ 为 $\Gamma$ 的长度。

\prop{2（可积函数必有界）}
设 $f(x,y,z)$ 在 $\Gamma$ 上可积，则其在 $\Gamma$ 上必有界。

\prop{3（积分的线性性质）}
设 $k_1,k_2$ 为常数，则
\[
\int_{\Gamma}[k_1f(x,y,z)\pm k_2g(x,y,z)]\dd s
=k_1\int_{\Gamma}f(x,y,z)\dd s
\pm k_2\int_{\Gamma}g(x,y,z)\dd s.
\]

\prop{4（积分的可加性）}
设 $f(x,y,z)$ 在 $\Gamma$ 上可积，且 $\Gamma_1\cup\Gamma_2=\Gamma$，$\Gamma_1\cap\Gamma_2=\varnothing$，则
\[
\int_{\Gamma}f(x,y,z)\dd s
=\int_{\Gamma_1}f(x,y,z)\dd s
+\int_{\Gamma_2}f(x,y,z)\dd s.
\]

\prop{5（积分的保号性）}
设 $f(x,y,z),g(x,y,z)$ 在 $\Gamma$ 上可积，且在 $\Gamma$ 上 $f(x,y,z)\le g(x,y,z)$，则
\[
\int_{\Gamma}f(x,y,z)\dd s\le\int_{\Gamma}g(x,y,z)\dd s.
\]
特殊地，有
\[
\left|\int_{\Gamma}f(x,y,z)\dd s\right|
\le\int_{\Gamma}|f(x,y,z)|\dd s.
\]

\prop{6（第一型曲线积分的估值定理）}
设 $M,m$ 分别是 $f(x,y,z)$ 在 $\Gamma$ 上的最大值和最小值，$l_{\Gamma}$ 为 $\Gamma$ 的长度，则
\[
ml_{\Gamma}\le\int_{\Gamma}f(x,y,z)\dd s\le Ml_{\Gamma}.
\]

\prop{7（第一型曲线积分的中值定理）}
设 $f(x,y,z)$ 在 $\Gamma$ 上连续，则在 $\Gamma$ 上至少存在一点 $(\xi,\eta,\zeta)$，使得
\[
\int_{\Gamma}f(x,y,z)\dd s=f(\xi,\eta,\zeta)l_{\Gamma}.
\]

\topictitle{3}{普通对称性与轮换对称性}

（1）普通对称性。假设 $\Gamma$ 关于 $xOz$ 面对称，$\Gamma_1$ 是 $\Gamma$ 在 $xOz$ 面右边的部分，则
\[
\int_{\Gamma}f(x,y,z)\dd s=
\begin{cases}
2\displaystyle\int_{\Gamma_1}f(x,y,z)\dd s,&f(x,y,z)=f(x,-y,z),\\[2mm]
0,&f(x,y,z)=-f(x,-y,z).
\end{cases}
\]
关于其他坐标面对称的情况与此类似。

（2）轮换对称性。若把 $x$ 与 $y$ 对调后，$\Gamma$ 不变，则
\[
\int_{\Gamma}f(x,y,z)\dd s=\int_{\Gamma}f(y,x,z)\dd s.
\]

\topictitle{4}{计算}

由于第一型曲线积分是由定积分推广而来的，因此计算第一型曲线积分的基本方法就是将其化为定积分。口诀为：\key{一投二代三计算}。

（1）平面情形。

若平面曲线 $L$ 由 $y=y(x)$（$a\le x\le b$）给出，则
\[
\dd s=\sqrt{1+[y'(x)]^2}\dd x,
\]
且
\[
\int_L f(x,y)\dd s
=\int_a^bf[x,y(x)]\sqrt{1+[y'(x)]^2}\dd x.
\]

\Needspace{8\baselineskip}
若平面曲线 $L$ 由参数式
\[
x=x(t),\qquad y=y(t),\qquad \alpha\le t\le\beta
\]
给出，则
\[
\dd s=\sqrt{[x'(t)]^2+[y'(t)]^2}\dd t,
\]
且
\[
\int_L f(x,y)\dd s
=\int_{\alpha}^{\beta}f[x(t),y(t)]
\sqrt{[x'(t)]^2+[y'(t)]^2}\dd t.
\]

若平面曲线 $L$ 由极坐标形式 $r=r(\theta)$（$\alpha\le\theta\le\beta$）给出，则
\[
\dd s=\sqrt{[r(\theta)]^2+[r'(\theta)]^2}\dd\theta,
\]
且
\[
\int_Lf(x,y)\dd s
=\int_{\alpha}^{\beta}f[r(\theta)\cos\theta,r(\theta)\sin\theta]
\sqrt{r^2+(r')^2}\dd\theta.
\]

（2）空间情形。若空间曲线 $\Gamma$ 由参数式
\[
x=x(t),\qquad y=y(t),\qquad z=z(t),\qquad \alpha\le t\le\beta
\]
给出，则
\[
\dd s=\sqrt{[x'(t)]^2+[y'(t)]^2+[z'(t)]^2}\dd t,
\]
且
\[
\int_{\Gamma}f(x,y,z)\dd s
=\int_{\alpha}^{\beta}f[x(t),y(t),z(t)]
\sqrt{[x'(t)]^2+[y'(t)]^2+[z'(t)]^2}\dd t.
\]

\begin{notebox}
曲线（曲面）积分中的变量定义在线（面）上，不具有独立性，因此“二代”时必须把曲线（曲面）方程代入被积函数。
\end{notebox}

\topictitle{5}{应用}

（1）空间光滑曲线 $\Gamma$ 的长度为 $l_{\Gamma}=\int_{\Gamma}\dd s$。

（2）若线密度为 $\rho(x,y,z)$，则质心为
\[
\bar x=\frac{\int_{\Gamma}x\rho\dd s}{\int_{\Gamma}\rho\dd s},\qquad
\bar y=\frac{\int_{\Gamma}y\rho\dd s}{\int_{\Gamma}\rho\dd s},\qquad
\bar z=\frac{\int_{\Gamma}z\rho\dd s}{\int_{\Gamma}\rho\dd s}.
\]
\begin{notebox}
（1）在考研范围内，重心就是质心。

（2）当密度 $\rho$ 为常数时，重心就成了形心。

（3）当密度 $\rho$ 为常数时，可反用形心公式
\[
\int_{\Gamma}x\dd s=\bar x\,l_{\Gamma},\qquad
\int_{\Gamma}y\dd s=\bar y\,l_{\Gamma},\qquad
\int_{\Gamma}z\dd s=\bar z\,l_{\Gamma}.
\]
\end{notebox}

（3）转动惯量为
\begin{align*}
I_x&=\int_{\Gamma}(y^2+z^2)\rho\dd s,&
I_y&=\int_{\Gamma}(x^2+z^2)\rho\dd s,\\
I_z&=\int_{\Gamma}(x^2+y^2)\rho\dd s,&
I_O&=\int_{\Gamma}(x^2+y^2+z^2)\rho\dd s.
\end{align*}

\parttitle{三}{第一型曲面积分}

\topictitle{1}{概念}

设曲面 $\Sigma$ 是光滑的，函数 $f(x,y,z)$ 在 $\Sigma$ 上有界。把 $\Sigma$ 任意分成 $n$ 个小块 $\Delta S_i$，设 $(\xi_i,\eta_i,\zeta_i)$ 是 $\Delta S_i$ 上任意取定的一点，作乘积 $f(\xi_i,\eta_i,\zeta_i)\Delta S_i$ 并作和。如果当各小块曲面的直径的最大值 $\lambda$ 趋于零时，该和的极限总存在，并且与 $\Delta S_i$ 的分法及取点的方法均无关，则称此极限为函数 $f(x,y,z)$ 在曲面 $\Sigma$ 上对面积的曲面积分或第一型曲面积分，记作
\[
\iintS f(x,y,z)\dd S
=\lim_{\lambda\to0}\sum_{i=1}^{n}f(\xi_i,\eta_i,\zeta_i)\Delta S_i.
\]

\begin{notebox}
（1）第一型曲面积分的物理背景是以 $f(x,y,z)$ 为面密度的空间物质曲面的质量。

（2）第一型曲面积分是由二重积分变换来的。
\end{notebox}

\topictitle{2}{性质}

以下均假设 $\Sigma$ 为空间有限分片光滑曲面。

\prop{1（求空间曲面的面积）}
\[
\iintS1\dd S=S,
\]
其中 $S$ 为 $\Sigma$ 的面积。

\prop{2（可积函数必有界）}
设 $f(x,y,z)$ 在 $\Sigma$ 上可积，则其在 $\Sigma$ 上必有界。

\prop{3（积分的线性性质）}
\[
\iintS[k_1f(x,y,z)\pm k_2g(x,y,z)]\dd S
=k_1\iintS f(x,y,z)\dd S
\pm k_2\iintS g(x,y,z)\dd S.
\]

\prop{4（积分的可加性）}
若 $\Sigma_1\cup\Sigma_2=\Sigma$，$\Sigma_1\cap\Sigma_2=\varnothing$，则
\[
\iintS f(x,y,z)\dd S
=\iint_{\Sigma_1}f(x,y,z)\dd S
+\iint_{\Sigma_2}f(x,y,z)\dd S.
\]

\prop{5（积分的保号性）}
若在 $\Sigma$ 上 $f(x,y,z)\le g(x,y,z)$，则
\[
\iintS f(x,y,z)\dd S\le\iintS g(x,y,z)\dd S.
\]
特殊地，有
\[
\left|\iintS f(x,y,z)\dd S\right|
\le\iintS|f(x,y,z)|\dd S.
\]

\prop{6（第一型曲面积分的估值定理）}
设 $M,m$ 分别是 $f(x,y,z)$ 在 $\Sigma$ 上的最大值和最小值，$S$ 为 $\Sigma$ 的面积，则
\[
mS\le\iintS f(x,y,z)\dd S\le MS.
\]

\prop{7（第一型曲面积分的中值定理）}
设 $f(x,y,z)$ 在 $\Sigma$ 上连续，则在 $\Sigma$ 上至少存在一点 $(\xi,\eta,\zeta)$，使得
\[
\iintS f(x,y,z)\dd S=f(\xi,\eta,\zeta)S.
\]

\topictitle{3}{普通对称性与轮换对称性}

（1）普通对称性。分析方法与二重积分、三重积分和第一型曲线积分完全一样。若 $\Sigma$ 关于 $xOz$ 面对称，$\Sigma_1$ 是 $\Sigma$ 在 $xOz$ 面右边的部分，则
\[
\iintS f(x,y,z)\dd S=
\begin{cases}
2\displaystyle\iint_{\Sigma_1}f(x,y,z)\dd S,&f(x,y,z)=f(x,-y,z),\\[2mm]
0,&f(x,y,z)=-f(x,-y,z).
\end{cases}
\]
（2）轮换对称性。若把 $x$ 与 $y$ 对调后 $\Sigma$ 不变，则
\[
\iintS f(x,y,z)\dd S=\iintS f(y,x,z)\dd S.
\]

\topictitle{4}{计算}

将曲面 $\Sigma$ 投影到某一坐标面时，要求 $\Sigma$ 上的任意两点在该坐标面上的投影点不重合；否则应把曲面分片。

若曲面写成 $z=z(x,y)$，其在 $xOy$ 面上的投影区域为 $D_{xy}$，则
\[
\dd S=\sqrt{1+z_x^2+z_y^2}\dd x\dd y,
\]
\[
\iintS f(x,y,z)\dd S
=\iint_{D_{xy}}f[x,y,z(x,y)]
\sqrt{1+z_x^2+z_y^2}\dd x\dd y.
\]
同理，
\[
\Sigma:y=y(x,z),\qquad
\dd S=\sqrt{1+y_x^2+y_z^2}\dd x\dd z,
\]
\[
\Sigma:x=x(y,z),\qquad
\dd S=\sqrt{1+x_y^2+x_z^2}\dd y\dd z.
\]

常用曲面面积元素：
\[
\begin{array}{c|c}
\text{曲面}&\text{面积元素}\\ \hline
x^2+y^2=R^2&\displaystyle\dd S=\frac{R}{\sqrt{R^2-x^2}}\dd x\dd z\\[5pt]
x^2+y^2+z^2=a^2&\displaystyle\dd S=\frac{a}{\sqrt{a^2-x^2-y^2}}\dd x\dd y\\[5pt]
z=\sqrt{x^2+y^2}&\displaystyle\dd S=\sqrt2\dd x\dd y
\end{array}
\]

\topictitle{5}{应用}

（1）曲面面积为 $S=\iintS\dd S$。

（2）若面密度为 $\rho(x,y,z)$，则质量为 $M=\iintS\rho(x,y,z)\dd S$，质心为
\[
\bar x=\frac{\iintS x\rho\dd S}{\iintS\rho\dd S},\qquad
\bar y=\frac{\iintS y\rho\dd S}{\iintS\rho\dd S},\qquad
\bar z=\frac{\iintS z\rho\dd S}{\iintS\rho\dd S}.
\]
\begin{notebox}
（1）在考研范围内，重心就是质心。

（2）当密度 $\rho$ 为常数时，重心就成了形心。

（3）当密度 $\rho$ 为常数时，可反用形心公式
\[
\iintS x\dd S=\bar xS,\qquad
\iintS y\dd S=\bar yS,\qquad
\iintS z\dd S=\bar zS.
\]
\end{notebox}

（3）转动惯量为
\begin{align*}
I_x&=\iintS(y^2+z^2)\rho\dd S,&
I_y&=\iintS(x^2+z^2)\rho\dd S,\\
I_z&=\iintS(x^2+y^2)\rho\dd S,&
I_O&=\iintS(x^2+y^2+z^2)\rho\dd S.
\end{align*}

\parttitle{四}{平面第二型曲线积分}

\topictitle{1}{变力沿曲线做功}

在一个向量场（变力场）中，设某质点在变力
\[
\bm F(x,y)=P(x,y)\bm i+Q(x,y)\bm j
\]
作用下，沿有向曲线 $L$ 从起点 $A$ 移动到终点 $B$。设在 $M(x,y)$ 点移动了微位移
\[
\dd\bm s=\dd x\bm i+\dd y\bm j,
\]
则力在此微位移上的微功为
\[
\dd W=\bm F(x,y)\cdot\dd\bm s.
\]
从而总功为
\[
W=\int_L\dd W
=\int_L\bm F(x,y)\cdot\dd\bm s
=\int_LP(x,y)\dd x+Q(x,y)\dd y.
\]

\topictitle{2}{概念}

第二型曲线积分的被积函数
\[
\bm F(x,y)=P(x,y)\bm i+Q(x,y)\bm j
\]
定义在平面有向曲线 $L$ 上，其物理背景是变力 $\bm F(x,y)$ 在平面曲线 $L$ 上从起点移动到终点所做的总功：
\[
\int_LP(x,y)\dd x+Q(x,y)\dd y.
\]

\topictitle{3}{性质}

以下总假设 $\Gamma$ 为有限长分段光滑有向曲线。

\prop{1（积分的线性性质）}
\[
\int_{\Gamma}(k_1\bm F_1\pm k_2\bm F_2)\cdot\dd\bm s
=k_1\int_{\Gamma}\bm F_1\cdot\dd\bm s
\pm k_2\int_{\Gamma}\bm F_2\cdot\dd\bm s.
\]

\prop{2（积分的有向性）}
\[
\int_{\widehat{AB}}\bm F\cdot\dd\bm s
=-\int_{\widehat{BA}}\bm F\cdot\dd\bm s.
\]

\prop{3（积分的可加性）}
当 $\widehat{AC}=\widehat{AB}+\widehat{BC}$ 时，
\[
\int_{\widehat{AC}}\bm F\cdot\dd\bm s
=\int_{\widehat{AB}}\bm F\cdot\dd\bm s
+\int_{\widehat{BC}}\bm F\cdot\dd\bm s.
\]

\begin{notebox}[第二型曲线积分的“对称性”]
若 $L$ 为从起点 $A$ 到终点 $B$ 的有向曲线，则它并不关于坐标轴对称。将对称两点处的 $\dd\bm s$ 作水平、垂直分解，两处的 $\dd x$ 同向，而 $\dd y$ 反向。若有力 $\bm F=x^2y\bm j$ 沿 $L$ 做功，则两处功的微元分别为 $x^2y\dd y$ 与 $-x^2y\dd y$，加起来为零，故 $\int_Lx^2y\dd y=0$。

这与数量积分完全不同：若 $f(x,y)=x^2y$，$L$ 为关于 $y$ 轴对称的曲线，则 $\int_Lf(x,y)\dd s=2\int_{L_1}x^2y\dd s$。数量积分的对称性是基于几何背景的；向量积分是物理问题的数学表达，多了“方向”这个重要因素。所以严格来讲，第二型曲线、曲面积分没有几何上的对称性，而是“在物理概念的基础上，得出数学表达式并确定好方向后，在数量大小上看是否相等”。
\end{notebox}

\topictitle{4}{计算}
\paragraph{（1）基本方法——化为定积分}

如果平面有向曲线 $L$ 由参数方程
\[
x=x(t),\qquad y=y(t),\qquad t:\alpha\to\beta
\]
给出，其中 $t=\alpha$ 对应起点 $A$，$t=\beta$ 对应终点 $B$，则
\[
\int_LP\dd x+Q\dd y
=\int_{\alpha}^{\beta}
\bigl\{P[x(t),y(t)]x'(t)+Q[x(t),y(t)]y'(t)\bigr\}\dd t.
\]

\begin{notebox}
$\alpha,\beta$ 的大小关系不作要求，关键是它们分别对应起点和终点。
\end{notebox}

\paragraph{（2）格林公式}

设 $D$ 为平面有界闭区域，$L$ 为 $D$ 的分段光滑边界，函数 $P(x,y),Q(x,y)$ 在包含 $D$ 的开区域内具有一阶连续偏导数。若 $L$ 取正向，则
\[
\boxed{
\ointL P\dd x+Q\dd y
=\iint_D\left(\frac{\partial Q}{\partial x}-\frac{\partial P}{\partial y}\right)\dd x\dd y
}.
\]
$L$ 的正向是指当一个人沿 $L$ 前进时，左手始终在 $L$ 所围成的区域 $D$ 内。

\circitem{1}\key{曲线封闭且无奇点在其内部，直接用格林公式。}
给定闭曲线积分 $\oint_LP\dd x+Q\dd y$ 时，检验 $P,Q$ 在该闭曲线所包围的区域内是否具有一阶连续偏导数；若满足，则直接使用格林公式。

\circitem{2}\key{非封闭曲线且 $Q_x\ne P_y$，补线使其封闭（加线减线）。}
补一条容易计算的曲线，使原曲线与补线构成一条封闭曲线。设所围区域为 $D$，则
\[
\int_LP\dd x+Q\dd y
=\pm\iint_D(Q_x-P_y)\dd x\dd y
-\int_{L_0}P\dd x+Q\dd y,
\]
其中正负号由闭合曲线的方向确定。

\circitem{3}\key{曲线封闭但有奇点在其内部，且除奇点外 $Q_x-P_y=0$，换路径。}
若两条闭曲线包围相同的奇点且方向相同，则可用一条容易计算的闭曲线替换原曲线。

格林公式给出的面积公式为
\[
|D|=\oint_Lx\dd y=-\oint_Ly\dd x
=\frac12\oint_L(x\dd y-y\dd x).
\]

\paragraphline{（3）平面曲线积分与路径无关}

\circitem{1}\key{概念}

设函数 $P(x,y),Q(x,y)$ 在区域 $G$ 内具有一阶连续偏导数。如果对 $G$ 内任意指定的两点 $A,B$ 以及连接 $A,B$ 的任意两条曲线 $L_1,L_2$，恒有
\[
\int_{L_1}P\dd x+Q\dd y
=\int_{L_2}P\dd x+Q\dd y,
\]
则称曲线积分 $\int_LP\dd x+Q\dd y$ 在 $G$ 内与路径无关。

曲线积分与路径无关相当于沿 $G$ 内任意闭曲线 $C$ 的曲线积分为零：
\[
\oint_CP\dd x+Q\dd y=0.
\]

\circitem{2}\key{条件}

设 $P,Q$ 在单连通区域 $G$ 内具有一阶连续偏导数，则曲线积分
$\int_LP\dd x+Q\dd y$ 在 $G$ 内与路径无关（或沿任意闭曲线的曲线积分为零）的充分必要条件是在 $G$ 内处处有
\[
\frac{\partial P}{\partial y}=\frac{\partial Q}{\partial x},
\]
或在 $G$ 内存在函数 $u(x,y)$，使得
\[
\dd u=P\dd x+Q\dd y.
\]

\begin{notebox}
（1）若平面区域内任一闭曲线所围的部分都属于该区域，则称该区域为平面单连通区域，否则称为复连通区域。通俗地说，平面单连通区域就是不含有“洞”（包含点“洞”）的区域，复连通区域是含有“洞”（包含点“洞”）的区域。

（2）设 $D$ 是单连通区域，且 $P,Q$ 具有一阶连续偏导数，则以下六个命题等价：

\circnum{1} $\dfrac{\partial Q}{\partial x}=\dfrac{\partial P}{\partial y}$（旋度为零的等式）在 $D$ 内处处成立；

\circnum{2} 沿 $D$ 内任意分段光滑闭曲线 $L$ 都有 $\displaystyle\oint_LP\dd x+Q\dd y=0$；

\circnum{3} $\displaystyle\int_{L_1}P\dd x+Q\dd y=\int_{L_2}P\dd x+Q\dd y$（积分与路径无关）；

\circnum{4} $\dd u=P\dd x+Q\dd y$（$P\dd x+Q\dd y$ 为某二元函数 $u(x,y)$ 的全微分）；

\circnum{5} $P\dd x+Q\dd y=0$ 是全微分方程；

\circnum{6} $(P,Q)$ 是某二元函数 $u$ 的梯度。
\end{notebox}

\circitem{3}\key{计算}

若 $u_x=P,u_y=Q$，则原函数可由折线路径求得：
\[
u(x,y)=\int_{x_0}^{x}P(t,y_0)\dd t
+\int_{y_0}^{y}Q(x,t)\dd t+C,
\]
或
\[
u(x,y)=\int_{y_0}^{y}Q(x_0,t)\dd t
+\int_{x_0}^{x}P(t,y)\dd t+C.
\]

\parttitle{五}{第二型曲面积分}

\topictitle{1}{向量场的通量}

如果 $\Omega$ 上的每一点 $M(x,y,z)$ 都对应着一个向量 $\bm F$，则在 $\Omega$ 上确定了向量函数
\[
\bm F(x,y,z)=P(x,y,z)\bm i+Q(x,y,z)\bm j+R(x,y,z)\bm k,
\]
它表示一个向量场。

在一个向量场中，设 $\Sigma$ 是某一有向分片光滑曲面，并指定了曲面的外侧，则向量函数 $\bm F(x,y,z)$ 通过曲面 $\Sigma$ 的通量为
\[
\iintS\bm F\cdot\dd\bm S
=\iintS\bm F\cdot\bn\dd S,
\]
其中
\[
\bn=(\cos\alpha,\cos\beta,\cos\gamma),
\qquad
\dd\bm S=(\dd y\dd z,\dd z\dd x,\dd x\dd y).
\]
因此
\[
\iintS\bm F\cdot\dd\bm S
=\iintS P(x,y,z)\dd y\dd z
+Q(x,y,z)\dd z\dd x
+R(x,y,z)\dd x\dd y.
\]

\topictitle{2}{概念}

第二型曲面积分的被积函数
\[
\bm F(x,y,z)=P(x,y,z)\bm i+Q(x,y,z)\bm j+R(x,y,z)\bm k
\]
定义在光滑的空间有向曲面 $\Sigma$ 上，其物理背景是向量函数 $\bm F(x,y,z)$ 通过曲面 $\Sigma$ 的通量：
\[
\iintS P\dd y\dd z+Q\dd z\dd x+R\dd x\dd y.
\]

\topictitle{3}{性质}

以下总假设 $\Sigma$ 是有向分片光滑曲面。

\prop{1（积分的线性性质）}
\[
\iintS(k_1\bm F_1\pm k_2\bm F_2)\cdot\dd\bm S
=k_1\iintS\bm F_1\cdot\dd\bm S
\pm k_2\iintS\bm F_2\cdot\dd\bm S.
\]

\prop{2（积分的方向性）}
\[
\iint_{\Sigma^-}\bm F\cdot\dd\bm S
=-\iint_{\Sigma^+}\bm F\cdot\dd\bm S.
\]

\prop{3（积分的可加性）}
当 $\Sigma_1\cup\Sigma_2=\Sigma$，$\Sigma_1\cap\Sigma_2=\varnothing$，且各曲面取向相容时，
\[
\iintS\bm F\cdot\dd\bm S
=\iint_{\Sigma_1}\bm F\cdot\dd\bm S
+\iint_{\Sigma_2}\bm F\cdot\dd\bm S.
\]

\begin{notebox}[第二型曲面积分的“对称性”]
与第二型曲线积分类似，第二型曲面积分没有几何上的对称性。设有向曲面 $\Sigma$ 关于 $xOz$ 面对称，$Q(x,y,z)=xy^2z$（关于 $y$ 的偶函数），则 $\iintS Q(x,y,z)\dd z\dd x=\iintS xy^2z\dd z\dd x=0$。

\circnum{1} 对称两处 $\dd\bm S$ 的法向量在 $\bm j$ 方向上的投影方向相反，故一处为 $Q(x,y,z)\dd z\dd x=xy^2z\dd z\dd x$，另一处为 $Q(x,-y,z)(-\dd z\dd x)=-xy^2z\dd z\dd x$，于是 $\iintS xy^2z\dd z\dd x=0$。

\circnum{2} 从通量的角度理解，一般规定流入为负通量、流出为正通量；若从一处流入、从对称处流出，通量为零，故积分为零。
\end{notebox}

\topictitle{4}{计算}
\paragraphline{（1）基本方法——化为二重积分}

\circitem{1}\key{拆成三个积分（如果有的话），一个一个做：}
\[
\iintS P\dd y\dd z+Q\dd z\dd x+R\dd x\dd y
=\iintS P\dd y\dd z+\iintS Q\dd z\dd x+\iintS R\dd x\dd y.
\]

\circitem{2}\key{分别投影到相应的坐标面上。}
例如，对于 $\iintS R(x,y,z)\dd x\dd y$，将曲面 $\Sigma$ 投影到 $xOy$ 平面上；$\dd x\dd y$ 只能往 $xOy$ 面投影，$\dd z\dd x$ 只能往 $xOz$ 面投影，$\dd y\dd z$ 只能往 $yOz$ 面投影。

\begin{notebox}
（1）对于第一型曲面积分而言，投影时并不允许随便选择投影方向，要确保曲面方程能写成相应单值函数；如果写不出，说明此时投影方向是错的。

（2）在第二型曲面积分中，$\dd y\dd z,\dd z\dd x,\dd x\dd y$ 不能想换什么换什么；一旦要转换投影，要先经过计算。转化为 $\dd x\dd y$ 才能往 $xOy$ 面投影，转化为 $\dd z\dd x$ 只能往 $xOz$ 面投影，转化为 $\dd y\dd z$ 只能往 $yOz$ 面投影。例如，对于圆柱面 $\Sigma$，若它在 $xOy$ 面上的投影为一条线，则 $\iintS R(x,y,z)\dd x\dd y=\iint_{D_{xy}}R\dd x\dd y=0$。
\end{notebox}

\circitem{3}\key{一投二代三计算。}
若曲面 $\Sigma:z=z(x,y)$，且在 $xOy$ 面上的投影区域为 $D_{xy}$，则一投，确定投影区域 $D_{xy}$；二代，将 $z=z(x,y)$ 代入被积函数；三计算，根据曲面的侧把 $\dd x\dd y$ 写成 $\pm\dd x\dd y$：
\[
\iintS P\dd y\dd z+Q\dd z\dd x+R\dd x\dd y
=\pm\iint_{D_{xy}}(-Pz_x-Qz_y+R)\dd x\dd y.
\]
当上侧为正时取“$+$”，下侧为正时取“$-$”（上正下负）。

\begin{notebox}
上式等号左边是第二型曲面积分，$\dd x\dd y$ 为有向曲面微元在 $xOy$ 平面上的投影分量；等号右边是 $xOy$ 平面上的二重积分，$\dd x\dd y$ 为二重积分的面积微元。两者写法一样，但意义不一样。
\end{notebox}

\circitem{4}\key{计算已转化成的二重积分。}

\paragraphline{（2）转换投影法}

\circitem{1}\key{转换投影法中有向曲面正向单位法向量的求法。}
若 $\Sigma:z=z(x,y)$，则
\[
\bn=\pm\frac{(-z_x,-z_y,1)}{\sqrt{1+z_x^2+z_y^2}},
\]
上侧为正取“$+$”，下侧为正取“$-$”（上正下负）。

\begin{notebox}
若 $\Sigma:y=y(x,z)$，则
\[
\bn=\pm\frac{(-y_x,1,-y_z)}{\sqrt{1+y_x^2+y_z^2}},
\]
右侧为正取“$+$”，左侧为正取“$-$”（右正左负）。

若 $\Sigma:x=x(y,z)$，则
\[
\bn=\pm\frac{(1,-x_y,-x_z)}{\sqrt{1+x_y^2+x_z^2}},
\]
前侧为正取“$+$”，后侧为正取“$-$”（前正后负）。
\end{notebox}

\circitem{2}\key{转换投影定理。}
设曲面 $\Sigma:z=z(x,y)$，$z$ 有一阶连续偏导数，且 $P,Q,R$ 在 $\Sigma$ 上连续，则
\[
\iintS P\dd y\dd z+Q\dd z\dd x+R\dd x\dd y
=\pm\iint_{D_{xy}}(-Pz_x-Qz_y+R)\dd x\dd y.
\]

第二型曲面积分与第一型曲面积分的关系为
\[
\iintS P\dd y\dd z+Q\dd z\dd x+R\dd x\dd y
=\iintS(P\cos\alpha+Q\cos\beta+R\cos\gamma)\dd S.
\]

\paragraph{（3）高斯公式}

设空间有界闭区域 $\Omega$ 由有向分片光滑闭曲面 $\Sigma$ 围成，$P(x,y,z),Q(x,y,z),R(x,y,z)$ 在 $\Omega$ 上具有一阶连续偏导数，则
\[
\boxed{
\oiintS P\dd y\dd z+Q\dd z\dd x+R\dd x\dd y
=\iiintO\left(
\frac{\partial P}{\partial x}
+\frac{\partial Q}{\partial y}
+\frac{\partial R}{\partial z}
\right)\dd v
},
\]
其中 $\Sigma$ 是 $\Omega$ 的整个边界曲面的外侧。

\Needspace{12\baselineskip}
\begin{notebox}
第二型曲面积分与“源”的概念是紧密联系的。空间区域内某点的散度是指这个点发散的强度；若区域内每一点都满足
\[
\frac{\partial P}{\partial x}
+\frac{\partial Q}{\partial y}
+\frac{\partial R}{\partial z}=0,
\]
则该场没有源头，称为“无源场”。
\end{notebox}

\circitem{1}\key{封闭曲面且内部无奇点，直接用高斯公式。}

\circitem{2}\key{非封闭曲面且 $\diver\bm F\ne0$，补面使其封闭（加面减面）。}

\circitem{3}\key{封闭曲面，有奇点在其内部，且除奇点外 $\diver\bm F=0$，可换一个面积分。}
若替换曲面与原曲面包围相同的奇点并取同向，则两个曲面积分相等；边界无须与原曲面重合。

若闭曲面 $\Sigma$ 包围原点并取外侧，则
\[
\oiintS
\frac{x\dd y\dd z+y\dd z\dd x+z\dd x\dd y}
{(x^2+y^2+z^2)^{3/2}}
=4\pi.
\]

\parttitle{六}{空间第二型曲线积分的计算}

\circitem{1}\key{一投二代三计算（基本方法）}

设空间有向曲线 $\Gamma$ 由
\[
x=x(t),\qquad y=y(t),\qquad z=z(t),\qquad t:\alpha\to\beta
\]
给出，则
\[
\int_{\Gamma}P\dd x+Q\dd y+R\dd z
=\int_{\alpha}^{\beta}
\Bigl\{P[x(t),y(t),z(t)]x'(t)
+Q[x(t),y(t),z(t)]y'(t)
+R[x(t),y(t),z(t)]z'(t)\Bigr\}\dd t.
\]

\circitem{2}\key{用斯托克斯公式}

设 $\Omega$ 为某空间区域，$\Sigma$ 为 $\Omega$ 内的分片光滑有向曲面片，$\Gamma$ 为逐段光滑的 $\Sigma$ 的边界；$\Gamma$ 的方向与 $\Sigma$ 的法向量成右手系。函数 $P(x,y,z),Q(x,y,z),R(x,y,z)$ 在 $\Omega$ 内具有连续的一阶偏导数，则
\[
\boxed{
\oint_{\Gamma}P\dd x+Q\dd y+R\dd z
=\iint_{\Sigma}
\begin{vmatrix}
\dd y\dd z&\dd z\dd x&\dd x\dd y\\
\dfrac{\partial}{\partial x}&\dfrac{\partial}{\partial y}&\dfrac{\partial}{\partial z}\\
P&Q&R
\end{vmatrix}
}.
\]
也可写成第一型曲面积分形式
\[
\oint_{\Gamma}P\dd x+Q\dd y+R\dd z
=\iint_{\Sigma}
\begin{vmatrix}
\cos\alpha&\cos\beta&\cos\gamma\\
\dfrac{\partial}{\partial x}&\dfrac{\partial}{\partial y}&\dfrac{\partial}{\partial z}\\
P&Q&R
\end{vmatrix}\dd S,
\]
其中 $\bn=(\cos\alpha,\cos\beta,\cos\gamma)$ 为 $\Sigma$ 的单位法向量。

\par\noindent\key{【注】}\quad
公式的成立与绷在 $\Gamma$ 上的曲面大小、形状无关。若 $\Sigma_1,\Sigma_2$ 具有相同边界 $\Gamma$，且方向与 $\Gamma$ 均按右手法则相配，则
\[
\oint_{\Gamma}=\iint_{\Sigma_1}=\iint_{\Sigma_2}.
\]
